\documentclass[11pt]{article}
\usepackage[margin=1in]{geometry}
\newtheorem{theorem}{Theorem}
\title{A Teaching Example for Verified Risk Reporting}
\author{}
\date{}
\begin{document}
\maketitle
\vspace{-2em}

\section{A release rule}
Fix a rule before calibration, with population risk $r$ and $m$ independent
Bernoulli calibration losses of mean $r$. Let $\tau$ be the risk budget and
let $\delta$ be the allowed false-release probability. Release the rule if
all $m$ losses are zero and $(1-\tau)^m\leq\delta$.

\begin{theorem}\label{thm:teaching-release}
For any fixed rule under this calibration design, the probability over
calibration of releasing it with $r>\tau$ is at most $\delta$.
\end{theorem}
If $r>\tau$, the probability of observing no loss is
$(1-r)^m < (1-\tau)^m \leq \delta$ whenever release is possible. If
$r\leq\tau$, a release cannot have risk strictly above the budget.

\section{Synthetic results}
The six records are synthetic teaching data. The displayed standard error
uses the Bernoulli plug-in formula for the joint event rate.

Our independent replay documented in the accompanying synthetic teaching ledger found 2 of 6 repetitions with an actual release and population risk above the budget. The joint violation frequency is 33.33 percent with a Monte Carlo standard error of 19.25 percentage points.

\section{Conclusion}
The example illustrates reporting the joint event with its repetition
denominator and stating the scope of the calibration guarantee.
\end{document}
